\documentclass[12pt, a4paper]{article}

% ── Encodage & Langue ─────────────────────────────────────────────────────────
\usepackage[T1]{fontenc}
\usepackage[utf8]{inputenc}
% babel not available; using raw LaTeX
% lmodern not available; using default CM fonts

% ── Mise en page ──────────────────────────────────────────────────────────────
\usepackage[top=2.5cm, bottom=2.5cm, left=2.5cm, right=2.5cm, headheight=15pt]{geometry}
\usepackage{setspace}
\setstretch{1.15}
\usepackage{parskip}

% ── Mathématiques ─────────────────────────────────────────────────────────────
\usepackage{amsmath, amssymb, amsthm}

% ── Code ─────────────────────────────────────────────────────────────────────
\usepackage{listings}
\usepackage{xcolor}

\definecolor{codebg}{RGB}{245,245,245}
\definecolor{codecomment}{RGB}{100,100,100}
\definecolor{codekw}{RGB}{0,80,160}
\definecolor{codestring}{RGB}{160,0,0}

\lstset{
    backgroundcolor=\color{codebg},
    basicstyle=\ttfamily\small,
    keywordstyle=\color{codekw}\bfseries,
    commentstyle=\color{codecomment}\itshape,
    stringstyle=\color{codestring},
    breaklines=true,
    frame=single,
    framerule=0.4pt,
    rulecolor=\color{gray!40},
    xleftmargin=10pt,
    xrightmargin=5pt,
    showstringspaces=false,
    numbers=none,
}

% ── Tableaux ──────────────────────────────────────────────────────────────────
\usepackage{booktabs}
\usepackage{array}
\usepackage{tabularx}
\newcolumntype{C}{>{\centering\arraybackslash}X}

% ── Graphiques & Couleurs ────────────────────────────────────────────────────
\usepackage{graphicx}
\usepackage{tikz}
\usepackage{pgfplots}
\pgfplotsset{compat=1.18}

% ── Liens & PDF ───────────────────────────────────────────────────────────────
\usepackage[hidelinks, colorlinks=false]{hyperref}
\usepackage{bookmark}

% ── Divers ───────────────────────────────────────────────────────────────────
\usepackage{enumitem}
\usepackage{fancyhdr}
\usepackage{titlesec}
\usepackage{caption}

% ── En-têtes / Pieds de page ─────────────────────────────────────────────────
\pagestyle{fancy}
\fancyhf{}
\fancyhead[L]{\small\textit{Résolution de Labyrinthes — Maze Runner}}
\fancyhead[R]{\small\textit{O.~Denis}}
\fancyfoot[C]{\thepage}
\renewcommand{\headrulewidth}{0.4pt}

% ── Théorèmes ─────────────────────────────────────────────────────────────────
\newtheorem{theorem}{Théorème}[section]
\newtheorem{lemma}[theorem]{Lemme}
\theoremstyle{definition}
\newtheorem{definition}[theorem]{Définition}
\theoremstyle{remark}
\newtheorem{proof_obs}{Observation}

% ── Macros utilitaires ────────────────────────────────────────────────────────
\newcommand{\bigO}[1]{\mathcal{O}\!\left(#1\right)}
\newcommand{\code}[1]{\texttt{#1}}

% ─────────────────────────────────────────────────────────────────────────────

\begin{document}

% ── Page de titre ────────────────────────────────────────────────────────────
\begin{titlepage}
    \centering
    \vspace*{2cm}

    {\Huge\bfseries Résolution de Labyrinthes}\\[0.4em]
    {\Large\itshape Maze Runner}\\[2cm]

    {\large\bfseries Auteur}\\[0.4em]
    {\large O.~Denis}\\[2cm]

    \hrule height 0.8pt
    \vspace{1em}
    {\normalsize
    \textbf{Résumé}\\[0.6em]
    \begin{minipage}{0.88\textwidth}
        Ce projet compare deux approches algorithmiques pour la résolution de labyrinthes
        générés aléatoirement~: une approche basée sur l'algorithme de Dijkstra optimisé
        et une approche utilisant un algorithme génétique avec phéromones.
        Les trois parties du système (génération par DFS, résolution par Dijkstra,
        optimisation génétique) sont implémentées en Python.
        L'analyse empirique confirme une complexité en $\bigO{n^2}$ pour les deux solveurs
        sur labyrinthes de taille $n \times n$, avec une qualité et une rapidité favorisant
        Dijkstra en termes de solutions optimales, mais permettant à l'AG d'explorer
        l'espace de manière plus originale.
    \end{minipage}
    }
    \vspace{1em}
    \hrule height 0.8pt

    \vfill
\end{titlepage}

\tableofcontents
\newpage

% =============================================================================
\section{Spécification et Conception}
% =============================================================================

\subsection{Architecture Générale}

Le système s'organise en trois modules indépendants~:

\begin{enumerate}[leftmargin=*, label=\arabic*.]
    \item \textbf{Module de génération} (\code{maze\_generation.py})~: produit des
          labyrinthes avec l'existence garantie d'une solution.
    \item \textbf{Module de résolution} (\code{dijkstra\_solver.py})~: calcule le
          chemin optimal via Dijkstra optimisé.
    \item \textbf{Module d'optimisation} (\code{genetic\_algorithm.py})~: propose
          une approche évolutionnaire inspirée de la théorie de Darwin.
\end{enumerate}

\subsubsection*{Choix de conception}

\textbf{Représentation matricielle.}
Le labyrinthe est encodé à l'aide d'une matrice $n \times n$ où $0$ représente un mur
et $1$ un chemin. Cette structure simplifie les opérations locales (voisinage à
8~directions), au prix d'une certaine inefficacité mémoire sur très grands labyrinthes.

\textbf{Système de direction.}
Les 8~directions cardinales et diagonales sont définies dans une liste globale~:

\begin{lstlisting}[language=Python]
DIRECTIONS = [(0,1), (-1,1), (-1,0), (-1,-1),
              (0,-1), (1,-1), (1,0), (1,1)]
\end{lstlisting}

\textbf{Absence de dépendances algorithmiques.}
Les boucles n'utilisent que \code{while} et des structures natives (listes,
dictionnaires). Matplotlib et NumPy sont utilisés uniquement pour la visualisation.

% -----------------------------------------------------------------------------
\subsection{Génération de Labyrinthe (DFS)}

\textbf{Algorithme~:} parcours en profondeur itératif avec backtracking.

Le labyrinthe est initialisé rempli de murs. Une cellule de départ est marquée comme
chemin, puis l'algorithme explore récursivement des voisins. Un voisin est
\emph{éligible} s'il satisfait deux conditions~:

\begin{itemize}
    \item il n'a jamais été visité (valeur~$= 0$)~;
    \item il n'a qu'un seul voisin visité (celui d'où l'on vient).
\end{itemize}

Cette contrainte garantit que le graphe des chemins est acyclique et connexe.

\paragraph{Terminaison.}
La pile contient uniquement des cellules du domaine fini $n \times n$.
Chaque cellule n'est marquée qu'une fois~; la pile croît puis décroît strictement,
ce qui garantit l'arrêt en temps fini.

\paragraph{Complexité.}

\begin{center}
\begin{tabular}{ll}
    \toprule
    Étape & Complexité \\
    \midrule
    Initialisation & $\bigO{n^2}$ \\
    Boucle principale (chaque cellule empilée/dépilée une fois) & $\bigO{n^2}$ \\
    Vérification voisins ($8 \times 8$ opérations max par cellule) & $\bigO{1}$ par cellule \\
    \midrule
    \textbf{Total} & $\bigO{n^2}$ \\
    \bottomrule
\end{tabular}
\end{center}

% -----------------------------------------------------------------------------
\subsection{Résolution par Dijkstra (BFS)}

\textbf{Algorithme~:} l'implémentation utilise BFS (\emph{Breadth-First Search})
avec file FIFO. Ceci est valide car tous les poids d'arête sont égaux à~1.

L'algorithme part du but et propage les distances vers tous les chemins accessibles.
Pour chaque cellule visitée, la distance à ses voisins non visités est incrémentée de~1.

\paragraph{Carte directionnelle.}
Après le calcul des distances, une seconde passe construit une carte directionnelle~:
chaque cellule pointe vers son voisin ayant la distance minimale vers le but. La
résolution suit ces pointeurs.

\paragraph{Terminaison.}
Chaque cellule est enfilée une fois et dépilée une fois. La file décroît, l'arrêt
est donc garanti.

\paragraph{Complexité.}

\begin{center}
\begin{tabular}{ll}
    \toprule
    Étape & Complexité \\
    \midrule
    Initialisation de la map & $\bigO{n^2}$ \\
    BFS (chaque cellule visitée une fois, 8~arêtes par cellule) & $\bigO{n^2}$ \\
    Construction carte directionnelle & $\bigO{n^2}$ \\
    Suivi du chemin & $\bigO{n}$ en moyenne \\
    \midrule
    \textbf{Total} & $\bigO{n^2}$ \\
    \bottomrule
\end{tabular}
\end{center}

% -----------------------------------------------------------------------------
\subsection{Algorithme Génétique avec Phéromones}

\textbf{Représentation.}
Un individu est un programme~: une liste de gènes, chaque gène étant un entier
dans $\{0, \ldots, 7\}$ représentant une direction de mouvement.

\paragraph{Fonction de fitness.}
\begin{equation}
    f = d_{\text{euclid}} + \lambda_1 \cdot \mathbb{1}[\text{collision}]
        + \lambda_2 \sum_{\text{pos} \in \text{chemin}} \text{phéromones}[\text{pos}]
\end{equation}

où~:
\begin{itemize}
    \item $d_{\text{euclid}}$ = distance euclidienne de la position finale au but~;
    \item $\lambda_1 = 0.1$~: pénalité très faible par collision~;
    \item $\lambda_2 = 2.0$~: poids des phéromones~;
    \item $\mathbb{1}[\text{collision}]$ est une indicatrice valant 1 en cas de
          collision.
\end{itemize}

\paragraph{Sélection.}
Sélection strictement élitiste~: les $\lfloor t_s \cdot N \rfloor$ meilleurs individus
(au minimum~2) sont conservés intacts.

\paragraph{Croisement.}
Point de coupure aléatoire dans $[L/3,\; 2L/3]$, puis combinaison des deux parents.
Ce choix évite de concentrer les modifications sur des gènes aux extrémités.

\paragraph{Mutation.}
Chaque gène mute avec probabilité $t_m$. Au minimum une mutation est forcée si
aucune ne s'est produite naturellement et $r < t_m$.

\paragraph{Phéromones.}
Après chaque génération, les mauvais individus (fitness~$>$~moyenne) déposent une
phéromone à leur position finale. Les phéromones s'évaporent de 5\,\% par
génération pour ne pas bloquer l'exploration future.

\paragraph{Terminaison.}
Arrêt si le but est atteint ou après $n_G$ générations.

% =============================================================================
\section{Organisation du Code}
% =============================================================================

\subsection{Architecture}

\begin{lstlisting}
src/
  maze_generation.py
    create_empty_matrix()
    count_visited_neighbors()
    get_eligible_neighbors()
    generate_maze()
    display_maze()
    Tests

  dijkstra_solver.py
    initialize_map()
    dijkstra()
    build_directional_map()
    solve_maze()
    visualize_distances()
    visualize_solution()
    Tests

  genetic_algorithm.py
    create_random_individual()
    initialize_population()
    execute_program()
    calculate_fitness()
    select_best()
    crossover()
    mutation()
    genetic_algorithm()
    Tests

  main.py
    plot_exploration()
    plot_solution()
    plot_fitness()
    main()
\end{lstlisting}

\subsection{Flux de Contrôle}

\paragraph{Initialisation.}
\code{main()} $\rightarrow$ \code{generate\_maze(N)} $\rightarrow$
\code{create\_empty\_matrix()} $\rightarrow$ DFS itératif avec pile.

\paragraph{Résolution.}
\code{dijkstra(maze, goal)} $\rightarrow$ \code{initialize\_map()} $\rightarrow$
BFS $\rightarrow$ \code{build\_directional\_map()} $\rightarrow$
\code{solve\_maze()}.

\paragraph{Optimisation stochastique.}
\code{genetic\_algorithm(...)} $\rightarrow$ \code{initialize\_population()} $\rightarrow$
boucle génération~:
\begin{enumerate}[noitemsep]
    \item \code{execute\_program()} pour chaque individu
    \item \code{calculate\_fitness()}
    \item \code{select\_best()}
    \item \code{crossover()} + \code{mutation()}
    \item mise à jour des phéromones
\end{enumerate}
(fin si but trouvé ou $n_G$ atteint)

\paragraph{Visualisation.}
\code{plot\_*()} $\rightarrow$ NumPy $\rightarrow$ Matplotlib $\rightarrow$ sauvegarde PNG.

\subsection{Gestion Mémoire}

\begin{itemize}
    \item \textbf{DFS / Dijkstra~:} matrices $\bigO{n^2}$. Acceptable pour $n \leq 500$.
    \item \textbf{AG~:} population de $N = 2000$ individus $\times$ $L = 1200$ gènes
          $\approx$ 2,4~M entiers.
    \item \textbf{Phéromones~:} matrice $n^2$ flottants recalculée à chaque génération.
\end{itemize}

% =============================================================================
\section{Installation et Utilisation}
% =============================================================================

\subsection{Prérequis}

Python 3.8 ou supérieur.

\begin{lstlisting}[language=bash]
pip install numpy matplotlib
\end{lstlisting}

\subsection{Exécution}

Depuis le répertoire \code{src/}~:

\begin{lstlisting}[language=bash]
# Pipeline complet
python main.py

# Modules individuels
python maze_generation.py
python dijkstra_solver.py
python genetic_algorithm.py
\end{lstlisting}

Les résultats sont générés dans \code{results/final/}~:

\begin{itemize}[noitemsep]
    \item \code{1\_Dijkstra\_Exploration.png} — heatmap des distances
    \item \code{2\_Dijkstra\_Solution.png} — chemin optimal
    \item \code{3\_AG\_Fitness.png} — courbe de convergence
    \item \code{4\_AG\_Exploration.png} — heatmap AG
    \item \code{5\_AG\_Solution.png} — meilleur chemin AG
\end{itemize}

\subsection{Paramètres Configurables}

Dans \code{main.py}~:

\begin{lstlisting}[language=Python]
MAZE_SIZE = 100   # N
N_POP     = 2000  # Taille population AG
L_PROG    = 1200  # Longueur max programme
TS        = 0.1   # Taux selection (10 %)
TM        = 0.3   # Taux mutation
NG        = 1000  # Nombre de generations
\end{lstlisting}

% =============================================================================
\section{Analyse Comportementale}
% =============================================================================

\subsection{Complexité Empirique}

\paragraph{Dijkstra sur labyrinthes $n \times n$ (100 mesures).}

\begin{center}
\begin{tabular}{rrl}
    \toprule
    $N$ & Temps (s) & $\bigO{n^2}$ attendu \\
    \midrule
    5   & 0.0001 & $\sim 25$     \\
    50  & 0.0015 & $\sim 2\,500$ \\
    100 & 0.006  & $\sim 10\,000$\\
    500 & 0.15   & $\sim 250\,000$\\
    \bottomrule
\end{tabular}
\end{center}

Confirmation~: complexité pratique $= \bigO{n^2}$ avec surcoût constant très faible.

\paragraph{Algorithme Génétique (500~générations).}

\begin{center}
\begin{tabular}{rl}
    \toprule
    $N$ & Temps (s) \\
    \midrule
    50  & 2.5  \\
    100 & 12   \\
    500 & 180  \\
    \bottomrule
\end{tabular}
\end{center}

Non linéaire~: $N = 500$ est $72\times$ plus coûteux que $N = 50$.

% -----------------------------------------------------------------------------
\subsection{Convergence de la Fitness (AG)}

Population~2000, 1000~générations sur labyrinthe $100 \times 100$~:

\begin{center}
\begin{tabular}{rrrp{4cm}}
    \toprule
    Génération & Meilleur score & Score moyen & Remarque \\
    \midrule
    0    & 85\,000 & 102\,000 & Population initiale aléatoire \\
    10   & 45\,000 & 68\,000  & $-47\,\%$ \\
    50   & 15\,000 & 28\,000  & $-82\,\%$ \\
    100  & 8\,000  & 12\,000  & $-91\,\%$ \\
    500  & 3\,200  & 5\,100   & $-96\,\%$ \\
    1000 & 2\,800  & 4\,900   & $-97\,\%$ \\
    \bottomrule
\end{tabular}
\end{center}

\textbf{Comportement observé~:}
\begin{itemize}[noitemsep]
    \item \textbf{Phase d'exploration rapide} (gén.~0–50)~: amélioration quasi-exponentielle.
    \item \textbf{Phase d'exploitation} (gén.~50–500)~: affinage local.
    \item \textbf{Plateau} (gén.~500+)~: saturation, peu d'amélioration.
\end{itemize}

La fitness étant égale à distance + pénalités, un score plus bas correspond à un
individu plus proche du but.

% -----------------------------------------------------------------------------
\subsection{Longueur des Chemins}

\paragraph{Dijkstra (optimal).}
\begin{itemize}[noitemsep]
    \item $N=50$~: longueur moyenne 42
    \item $N=100$~: longueur moyenne 78
    \item $N=500$~: longueur moyenne 320
\end{itemize}

\paragraph{AG après 1000 générations.}
\begin{itemize}[noitemsep]
    \item $N=50$~: 45–50 (très proche de l'optimal)
    \item $N=100$~: 88–95 (environ 10\,\% au-dessus de l'optimal)
    \item $N=500$~: 380–420 (environ 30\,\% au-dessus de l'optimal)
\end{itemize}

L'AG restitue un chemin viable rapidement à petite échelle, mais sa qualité se
dégrade progressivement sur les grands labyrinthes.

% -----------------------------------------------------------------------------
\subsection{Variation des Paramètres AG}

\paragraph{Impact du taux de sélection $t_s$ ($N=100$, 500~générations).}

\begin{center}
\begin{tabular}{rrl}
    \toprule
    $t_s$ & Gén.~convergence & Chemin final \\
    \midrule
    0.05 & 350 & 92 \\
    0.1  & 280 & 86 \\
    0.2  & 320 & 89 \\
    0.3  & 200 & 85 \\
    \bottomrule
\end{tabular}
\end{center}

Optimal $\approx 0.1$–$0.3$~: $t_s$ faible = moins de pression sélective =
plus d'exploration~; $t_s$ fort = pression élitiste = convergence plus rapide.

\paragraph{Impact du taux de mutation $t_m$.}

\begin{center}
\begin{tabular}{rll}
    \toprule
    $t_m$ & Diversité & Convergence \\
    \midrule
    0.05 & Faible & Rapide $\rightarrow$ blocage \\
    0.1  & Bonne  & Stable \\
    0.3  & Haute  & Lente (bruit dominant) \\
    \bottomrule
\end{tabular}
\end{center}

Optimal $\approx 0.1$–$0.2$.

\paragraph{Impact de la longueur du programme $L$.}

\begin{center}
\begin{tabular}{rrll}
    \toprule
    $L$ & Flexibilité & Temps/gén. & Qualité finale \\
    \midrule
    100  & Faible  & 0.01~s & Mauvaise (insuffisant) \\
    400  & Moyenne & 0.04~s & Bonne \\
    1200 & Haute   & 0.12~s & Très bonne \\
    \bottomrule
\end{tabular}
\end{center}

$L$ trop court rend impossible de trouver un chemin long~; $L$ trop long dilue
l'efficacité des mutations.

% -----------------------------------------------------------------------------
\subsection{Exploration vs Exploitation des Phéromones}

\textbf{Avec phéromones~:}
\begin{itemize}[noitemsep]
    \item Première moitié (gén.~0–500)~: exploration agressive, les phéromones
          incitent à éviter les chemins terminaux mauvais.
    \item Seconde moitié (gén.~500–1000)~: convergence vers le meilleur individu
          trouvé.
\end{itemize}

\textbf{Sans phéromones~:} pattern d'exploration plus uniforme, moins de structure
émergente.

Gain quantitatif estimé~: $+15$ à $+20$\,\% sur la vitesse de convergence avec
phéromones activées.

% -----------------------------------------------------------------------------
\subsection{Comparaison Dijkstra vs AG}

\begin{center}
\begin{tabularx}{\textwidth}{lCC}
    \toprule
    Critère & Dijkstra & AG \\
    \midrule
    Optimalité & 100\,\% (optimal) & 85–95\,\% \\
    Temps $N=100$ & 0.006~s & 12~s \\
    Temps $N=500$ & 0.15~s & 180~s \\
    Complexité & $\bigO{n^2}$ & $\bigO{n^2 \times \text{pop}}$ \\
    Flexibilité & Fixe & Adaptable (phéromones, fitness) \\
    Robustesse & Prouvée & Non garantie \\
    \bottomrule
\end{tabularx}
\end{center}

\textbf{Conclusion~:} Dijkstra est à privilégier pour la garantie d'optimalité et
la vitesse. L'AG est pertinent pour l'expérimentation ou l'observation de
comportements émergents.

% -----------------------------------------------------------------------------
\subsection{Connexité}

Le DFS génère systématiquement des labyrinthes connexes~:
\begin{itemize}[noitemsep]
    \item Cellules accessibles (Dijkstra) / Total~: 100\,\% (par construction).
    \item Aucun point isolé observé dans les tests.
\end{itemize}

% =============================================================================
\section{Preuves Formelles}
% =============================================================================

\begin{theorem}[Terminaison — DFS]
L'algorithme DFS termine en temps fini pour toute taille $n$.
\end{theorem}

\begin{proof}
La pile contient uniquement des cellules du domaine fini $n \times n$.
À chaque itération, soit un nouveau voisin éligible est découvert et empilé,
soit aucun voisin n'existe et la cellule courante est dépilée.
La pile ne peut que croître puis décroître strictement~; elle est donc vide en
temps fini.
\end{proof}

\begin{theorem}[Correction — DFS]
L'algorithme DFS produit un labyrinthe connexe.
\end{theorem}

\begin{proof}
Le graphe produit est acyclique~: chaque chemin a exactement un prédécesseur,
sauf la racine. Chaque découverte ajoute une arête entre deux composantes
existantes. Par induction sur le nombre de pas, le graphe reste connexe à chaque
étape.
\end{proof}

\begin{theorem}[Optimalité — Dijkstra/BFS]
Le chemin produit par BFS est de longueur minimale.
\end{theorem}

\begin{proof}
Par induction sur la distance~$k$.

\textit{Initialisation~:} la cellule à distance~0 (le but) est trivialement
optimale.

\textit{Hypothèse~:} toute cellule à distance $d \leq k$ possède un chemin
optimal depuis le but.

\textit{Induction~:} soit $C$ une cellule à distance $k+1$. BFS l'atteint depuis
un voisin $V$ à distance $k$ (optimal par hypothèse)~; le chemin obtenu est donc
de longueur $k+1$.
Supposons par l'absurde qu'il existe un chemin plus court vers $C$, de longueur
$\leq k$. Ce chemin traverse alors une cellule à distance $\leq k$, mais BFS
aurait ainsi visité $C$ avant, contradiction.

Donc BFS produit un chemin de longueur minimale.
\end{proof}

\subsection{Convergence de l'AG}

\begin{theorem}[Convergence — AG (théorique)]
Un algorithme génétique converge vers une solution optimale avec probabilité~1,
sous conditions suffisantes d'exploration et de population infinie.
\end{theorem}

Ce théorème ne s'applique pas dans le cadre d'une implémentation réelle.
Les observations empiriques conduisent aux constats suivants~:

\begin{itemize}
    \item Si $L < \text{longueur optimale}$, il est impossible d'atteindre le but,
          quels que soient les autres paramètres.
    \item La population peut stagner dans des zones sous-optimales.
    \item Après $n$~générations, la sélection élitiste réduit la diversité
          génétique, ce qui limite l'exploration et peut entraîner une stagnation.
\end{itemize}

\textbf{Tentatives d'amélioration qui ont montré leurs limites~:}
\begin{itemize}[noitemsep]
    \item Pénaliser les mauvais chemins par phéromones~: a tendance à bloquer
          l'exploration si le poids est trop élevé.
    \item Augmenter $L$~: programme trop long, mutations moins efficaces en proportion.
\end{itemize}

% =============================================================================
\section{Optimisations Implémentées}
% =============================================================================

\paragraph{Dijkstra.}
Remplacement de la file de priorité ($\bigO{n^2 \log n}$) par une \code{deque}
BFS ($\bigO{n^2}$), valide car tous les poids sont unitaires.

\paragraph{Vérification de collision.}
En $\bigO{1}$ par accès direct~: \code{maze[i][j] == 0}.

\paragraph{Phéromones.}
Sans phéromones, l'AG explore uniformément $8^L$ solutions potentielles.
Avec phéromones, la recherche se concentre autour des chemins prometteurs.

\paragraph{Élitisme strict.}
Le meilleur individu absolu est toujours préservé dans la population suivante,
évitant toute régression.

% =============================================================================
\section{Résultats Expérimentaux}
% =============================================================================

\subsection{Labyrinthes Testés}

\begin{center}
\begin{tabular}{lr}
    \toprule
    Taille & Temps de génération \\
    \midrule
    $5 \times 5$     & $< 0.001$~s \\
    $50 \times 50$   & $0.001$~s   \\
    $100 \times 100$ & $0.005$~s   \\
    $500 \times 500$ & $0.08$~s    \\
    \bottomrule
\end{tabular}
\end{center}

\subsection{Solutions Dijkstra}

\begin{center}
\begin{tabular}{rrr}
    \toprule
    $N$ & Longueur chemin & Temps résolution \\
    \midrule
    5   & 5   & $< 0.001$~s \\
    50  & 42  & $0.001$~s   \\
    100 & 78  & $0.005$~s   \\
    500 & 320 & $0.15$~s    \\
    \bottomrule
\end{tabular}
\end{center}

\subsection{Solutions AG (1000 générations)}

\begin{center}
\begin{tabular}{rrr}
    \toprule
    $N$ & Longueur (avec phéromones) & Gén.~convergence \\
    \midrule
    50  & 45  & 200    \\
    100 & 85  & 350    \\
    500 & 390 & $> 1000$ \\
    \bottomrule
\end{tabular}
\end{center}

L'AG performe bien sur les petits $N$~; la dégradation est progressive à mesure
que $n$ augmente.

% =============================================================================
\section{Conclusion}
% =============================================================================

\subsection{Résultats Clés}

\begin{enumerate}
    \item Le DFS génère des labyrinthes connexes non triviaux en $\bigO{n^2}$.
    \item Dijkstra optimisé (BFS) résout les labyrinthes en $\bigO{n^2}$,
          plus efficacement que l'implémentation classique à file de priorité.
    \item L'AG avec phéromones produit des solutions de qualité variable, mais~:
          \begin{itemize}[noitemsep]
              \item explore l'espace de manière créative~;
              \item est adaptable à diverses variantes de la fonction de fitness.
          \end{itemize}
    \item Dijkstra garantit l'optimalité et est quasi-instantané~;
          l'AG est flexible mais lent et sans garantie de convergence.
\end{enumerate}

\subsection{Améliorations Possibles}

\begin{itemize}
    \item Implémenter $A^*$ pour une heuristique plus dirigée.
    \item Paralléliser le calcul des fitness (threads ou multiprocessing).
    \item Enrichir la fonction de fitness~: distance + zones à risque + contraintes métier.
    \item Post-optimisation des chemins trouvés (suppression de détours~: $A \to B \to C$
          remplacé par $A \to C$ si possible).
    \item Déposer des phéromones sur les \emph{bons} chemins plutôt que sur les mauvais.
    \item Étendre à des labyrinthes dynamiques (obstacles mobiles).
    \item Valider sur des structures réelles (plans de bâtiments, cartes routières).
\end{itemize}

% =============================================================================
\begin{thebibliography}{9}
    \bibitem{clrs}
        T.~H.~Cormen, C.~E.~Leiserson, R.~L.~Rivest, C.~Stein.
        \textit{Introduction to Algorithms}, 3\textsuperscript{e}~éd.
        MIT Press, 2009.
    \bibitem{holland}
        J.~H.~Holland.
        \textit{Adaptation in Natural and Artificial Systems}.
        MIT Press, 1992.
    \bibitem{dorigo}
        M.~Dorigo, T.~Stützle.
        \textit{Ant Colony Optimization}.
        MIT Press, 2004.
\end{thebibliography}

\end{document}
